\section{Introduction}
\label{sec:introduction}

Research agents can now search literature, edit repositories, run experiments, and solve technical problems. A correct-looking result, however, can rest on a fabricated citation, weakened test, leaked evaluation set, invalid statistical unit, or missing proof condition. In these settings, output quality alone is an incomplete measure of help: a system must make a useful substantive judgment and recognize defects that would invalidate that judgment.

Capability benchmarks increasingly test realistic research work. \textsc{PaperBench} measures paper replication with author-backed rubrics \cite{starace2025paperbench}; \textsc{RE-Bench} and \textsc{MLE-bench} evaluate open-ended machine-learning engineering \cite{wijk2024rebench,chan2025mlebench}; and \textsc{SWE-bench} uses executable tests for repository issues \cite{jimenez2024swebench}. Work on literature synthesis and mathematics adds expert-grounded and verifiable tasks \cite{wang2024trialmind,glazer2024frontiermath}. These benchmarks establish that capability is growing, but evaluator exploitation and weak tests show why high scores need not imply valid research behavior.

Safety evaluations approach the problem from the other direction. Sabotage evaluations condition risk on a model, scaffold, access policy, monitor, and audit budget \cite{benton2024sabotage}. Yet capability and concealed-failure recognition are seldom scored on the same instance with clean controls. This leaves a practical gap: we cannot tell whether a model that reaches the right research conclusion also notices the fault that makes an artifact unsafe to trust.

We address this gap with \audit, a paired capability--integrity evaluation summarized in \figref{fig:pipeline}. We author 32 balanced cases across four research domains. Each case has an independently keyed substantive choice and integrity choice; half are clean and half contain a consequential planted fault. We compare an ordinary expert-review prompt with a safety-aware prompt that explicitly separates solving from auditing while keeping the model, case, decoding, and output schema fixed.

\begin{figure}[t]
\centering
\setlength{\fboxsep}{7pt}
\fbox{\parbox{0.25\linewidth}{\centering \textbf{32 cases}\\4 research domains\\16 clean / 16 faulted}}
\hfill$\longrightarrow$\hfill
\fbox{\parbox{0.25\linewidth}{\centering \textbf{Paired review}\\ordinary vs. safety-aware\\3 Qwen3 sizes}}
\hfill$\longrightarrow$\hfill
\fbox{\parbox{0.25\linewidth}{\centering \textbf{Exact audit}\\substance + integrity\\paired tests + risk bounds}}
\caption{Overview of \audit. Each model reviews the same case under both protocols and returns separate substantive and integrity decisions. Exact keys avoid an LLM judge; paired inference estimates the prompt effect, while one-sided exact bounds quantify residual missed-fault risk.}
\label{fig:pipeline}
\end{figure}


Across 192 trials, the models answer the substantive question correctly in 96.4\% of cases. The safety-aware prompt improves both substantive and integrity accuracy by 3.1 percentage points, with no reverse discordances, but adjusted McNemar tests give $p=.50$. More importantly, the protocol misses 3/48 faults: the 6.25\% point estimate has a 15.4\% one-sided 95\% upper bound and fails our preregistered 10\% residual-risk criterion. Thus the pilot demonstrates bounded usefulness, not a safety certificate.

Our contributions are:
\begin{itemize}[leftmargin=*,itemsep=0pt,topsep=0pt]
    \item We propose a same-instance audit that separates substantive correctness, integrity recognition, joint success, and uncertainty.
    \item We release a balanced 32-case benchmark with exact labels across literature review, coding, experiment design, and mathematics.
    \item We conduct 192 real-model trials and a paired prompt ablation with case-clustered intervals and exact rare-event bounds.
    \item We show that high capability and zero clean-case false alarms can coexist with insufficient evidence for low residual risk.
\end{itemize}

\Secref{sec:related} positions the audit against capability and safety evaluations. \Secref{sec:method} describes the benchmark, protocols, and preregistered analysis. \Secref{sec:results} reports the outcomes, and \secref{sec:discussion} discusses their scope and implications.
